% !PRINTED 660
\weberpage{660}

%<<CONT
\weq{3}{\begin{aligned}
\eta &= c + \varrho\eta_1, \\
\eta_1 &= c_1 + \varrho\eta_2, \\
&\;\cdots \\
\eta_{r-1} &= c_{r-1} + \varrho\eta_r,
\end{aligned}}

from which it follows:

\weq{4}{\eta = c + c_1\varrho + c_2\varrho^2 + \cdots + c_{r-1}\varrho^{r-1} + \eta_r\varrho^r.}

The sequence of the constants $c, c_1, c_2, \ldots, c_{r-1}$ is
completely determined by $\eta$, and $\eta_1, \eta_2, \ldots, \eta_r$
is a sequence of functions in $\Omega$ whose denominator is not
divisible by $\pi$.

One might call this expression the \gesperrt{Taylor development} of the function
$\eta$ in powers of $\varrho$. A similar development can also be given
for a function $\eta$ whose denominator is divisible by $\pi$, only
that negative powers then occur as well: one can namely determine a
power $\varrho^s$ in such a way that the denominator of
$\varrho^s\eta$ is no longer divisible by $\pi$, and if one applies
the development (4) to this product, there follows:

\weq{5}{\eta = c\varrho^{-s} + c_1\varrho^{-s+1} + \cdots + c_{r-1}\varrho^{r-s-1} + \eta_r\varrho^{r-s}.}

\begin{center}
§ 186. \textbf{\gesperrt{Birational Transformation.}}
\end{center}

We now come to a subject which distinguishes the region of algebraic
functions above all else from that of algebraic numbers, and the
reason for this is that the theory of integral functions and
functionals has for functions by no means the character of invariance
that it has for numbers; this is birational transformation.

If $z_1$ is any non-constant function of the field $\Omega$, and
$N(t - z_1)$ is the $f$th power of an irreducible function of degree
$e$ (§ 172, 6.), then there subsists between $z$ and $z_1$ an
equation which with respect to $z_1$ is of degree $e$, and whose
degree with respect to $z$, once denominators and superfluous factors
have been cleared away, shall be denoted by $e_1$. Let this equation
be

\weq{1}{G\bigl(\overset{e}{z_1},\, \overset{e_1}{z}\bigr) = 0.}

Among the various equations of this form there is only one
irreducible one, and this is of the lowest possible degree both with
respect to $z$ and with respect to $z_1$ (Bd. I, § 20).

If $\theta$ is a primitive function of the field $\Omega$, then by
§ 172, (12)

\weq{2}{z_1^h\theta^k \qquad
\begin{aligned}
h &= 0, 1, \ldots, e-1 \\
k &= 0, 1, \ldots, f-1
\end{aligned}}
is
%CONT>>
